\begin{recipe}{Pasta met chorizo en andijvie}
\kicker[Hoofdgerecht,Vlees,Spaans]{Hoofdgerecht · vlees · Spaans}
\index[register]{Pasta met chorizo en andijvie}
\index[register]{Chorizo}
\index[register]{Andijvie}
\index[register]{Fusilli}

\meta{20 minuten \dvd Spaans}

\lettrine{S}{nelle} doordeweekse pasta waarbij de chorizo genoeg vet loslaat
om alles in te bakken zonder extra olie. De andijvie slinkt snel mee met de
slagroom en de zongedroogde tomaten gaan er koud bovenop.

\blockrule{Ingrediënten}
\begin{ingredients}
  \ing{300 g}{fusilli}\index[register]{Fusilli}
  \ing{250 g}{chorizoworst aan een stuk}\index[register]{Chorizo}
  \ing{1}{ui}
  \ing{1 zak}{gesneden andijvie}\index[register]{Andijvie}
  \ing{125 ml}{verse slagroom}
  \ing{150 g}{semi-zongedroogde tomaten}\index[register]{Zongedroogde tomaten}
  \ingb{versgemalen peper en zout}
\end{ingredients}

\blockrule{Bereiding}
\tip{Verhit de hapjespan droog: de chorizo bakt in zijn eigen vet en heeft geen extra olie nodig.}
\begin{steps}
  \item Kook de fusilli beetgaar volgens de aanwijzingen op de verpakking.
  \item Verwijder ondertussen het vel van de chorizo en snijd in blokjes.
        Snipper de ui.
  \item Verhit een droge hapjespan op middelhoog vuur. Bak de chorizo
        4 minuten. Voeg de ui toe en bak 3 minuten mee.
  \item Voeg de slagroom en andijvie toe en laat al omscheppend slinken.
  \item Giet de pasta af en schep door het andijviemengsel. Breng op smaak
        met peper en zout.
  \item Verdeel over borden en leg de zongedroogde tomaten erop.
\end{steps}
\end{recipe}
